\section{Virtual Campus and 360$^{\circ}$ Tour Engine}
\label{sec:tour}
\begin{figure}[!t]
\centering
\begin{tikzpicture}[
  font=\scriptsize, node distance=8mm,
  n/.style={draw, circle, minimum size=6mm, inner sep=0pt, fill=black!5},
  chain/.style={-{Latex[length=1.6mm]}, thick},
  xref/.style={-{Latex[length=1.6mm]}, densely dashed, gray},
  cf/.style={-{Latex[length=1.6mm]}, dash dot, black!60}
]
% ground floor chain
\node[n] (g1) {$g_1$};
\node[n, right=of g1] (g2) {$g_2$};
\node[n, right=of g2] (g3) {$g_3$};
\node[n, right=of g3] (g4) {$g_4$};
\node[left=1mm of g1] {head};
\node[right=1mm of g4] {tail};
\draw[chain, transform canvas={yshift=1pt}] (g1)--(g2);
\draw[chain, transform canvas={yshift=-1pt}] (g2)--(g1);
\draw[chain, transform canvas={yshift=1pt}] (g2)--(g3);
\draw[chain, transform canvas={yshift=-1pt}] (g3)--(g2);
\draw[chain, transform canvas={yshift=1pt}] (g3)--(g4);
\draw[chain, transform canvas={yshift=-1pt}] (g4)--(g3);
\node[above=5mm of g2, font=\scriptsize\itshape] {Ground Floor list (\code{previous}$\leftrightarrows$\code{next})};

% first floor chain
\node[n, below=16mm of g1] (f1) {$f_1$};
\node[n, right=of f1] (f2) {$f_2$};
\node[n, right=of f2] (f3) {$f_3$};
\node[left=1mm of f1] {head};
\draw[chain, transform canvas={yshift=1pt}] (f1)--(f2);
\draw[chain, transform canvas={yshift=-1pt}] (f2)--(f1);
\draw[chain, transform canvas={yshift=1pt}] (f2)--(f3);
\draw[chain, transform canvas={yshift=-1pt}] (f3)--(f2);
\node[below=5mm of f2, font=\scriptsize\itshape] {First Floor list};

% cross-floor stitch tail(g)->head(f)
\draw[cf, bend right=25] (g4) to node[right, font=\scriptsize] {stitch: \code{tail.next}$=f_1$} (f1);
\draw[cf, bend left=25] (f1) to node[left, font=\scriptsize] {\code{head.previous}$=g_4$} (g4);

% cross references
\node[n, above=8mm of g3, fill=black!12] (up) {$u$};
\draw[xref] (g3) to node[right, font=\scriptsize] {\code{crossReferences.floorUp}} (up);
\node[n, below right=5mm and 3mm of g4, fill=black!12] (room) {$r$};
\draw[xref] (g4) to node[below, font=\scriptsize] {\code{roomEntry}} (room);
\end{tikzpicture}
\caption{Panorama scene graph. Each floor is a doubly linked list (solid);
branch and vertical connections are separate cross-reference pointers
(dashed); a construction pass stitches each floor's tail to the next
floor's head (dash-dot), so sequential traversal crosses floors with no
boundary-aware code.}
\label{fig:scenegraph}
\end{figure}

\begin{figure}[!t]
\centering
\begin{tikzpicture}[
  font=\scriptsize, node distance=7mm and 12mm,
  s/.style={draw, rounded corners=3pt, minimum height=6mm, inner sep=3pt, fill=black!4},
  pf/.style={-{Latex[length=1.6mm]}}
]
\node[s] (idle) {idle};
\node[s, right=of idle] (play) {playing};
\node[s, above=of play] (pause) {paused};
\node[s, right=of play] (done) {completed};
\node[s, below=of play] (err) {error};
\node[s, below=of idle] (stop) {stopped};

\draw[pf] (idle) -- node[above]{\code{start()} (seed usable)} (play);
\draw[pf] (idle) to[bend right=15] node[left]{seed unusable} (err);
\draw[pf] (play) to[bend left=12] node[right]{\code{pause()}} (pause);
\draw[pf] (pause) to[bend left=12] node[left]{\code{resume()} (from phase $i$)} (play);
\draw[pf] (play) -- node[above]{no \code{node.next}} (done);
\draw[pf] (play) -- node[right]{next unusable / revisited} (err);
\draw[pf] (play) to[bend left=10] node[below]{\code{stop()} (token${+}{+}$)} (stop);
\draw[pf] (pause) to[bend left=30] (stop);
\draw[pf] (stop) to[bend left=20] node[below]{\code{restart()} $\to$ seed node} (play);
\draw[pf] (done) to[bend right=30] node[above]{\code{restart()}} (play);
\end{tikzpicture}
\caption{Guided-tour status automaton (\code{useGuidedTour}). A monotonic
cancellation token orphans in-flight timer chains on every transition out
of \emph{playing}; a per-run visited set aborts to \emph{error} on a
repeated scene.}
\label{fig:guided}
\end{figure}


The tour is not merely a set of 360$^{\circ}$ images: it is a data
structure over which two traversal modes operate. This section describes
the scene graph, the cross-floor stitching pass, the manual and guided
traversal algorithms, and the minimap projection, all as implemented in
\code{frontend/src/features/tour/engine/}.

\subsection{Scene-Graph Data Structure}
The overall structure is shown in Fig.~\ref{fig:scenegraph}. Each panorama
scene is a \code{PanoramaNode} object holding an image path,
a preview path, a resting orientation $(\mathrm{yaw}, \mathrm{pitch},
\mathrm{hfov})$, a $1$-based per-floor sequence index, \code{previous} and
\code{next} links, a \code{crossReferences} record, and a list of
hand-authored \code{crossFloorHotspots}. A link
$\ell = (\text{node}, \mathrm{entryYaw}, \mathrm{entryPitch})$ carries the
camera heading the target should open with when reached \emph{via that
link}, in the Street-View style~\cite{anguelov2010streetview} of preserving
the direction of travel.

Each floor is a \code{PanoramaLinkedList}: a doubly linked list with
\code{head}/\code{tail} pointers and a \code{Map}$\langle$id$,$node$\rangle$
for $O(1)$ random access. Sequential traversal is $O(1)$ per step; the
list's \code{toArray()} deliberately stops at its own \code{tail} rather
than following \code{next} indefinitely, because after stitching a tail's
\code{next} points into another floor. The whole building is a
\code{PanoramaEngine}: a \code{Map}$\langle$floorName$,$list$\rangle$ plus a
flat node index, built once per tour-page render by \code{useMemo} directly
from the backend \code{/tour/scenes} response --- never a hardcoded array.

\emph{Why a doubly linked list.} A flat array with \code{findIndex} (the
prior design) makes ``next scene'' an $O(n)$ scan and couples every
consumer to array positions; a modulo wraparound over the array also
produces an artificial jump from the last scene back to the first. A
general graph structure would support the branch connections but makes
``the next scene in the walking sequence'' ill-defined. The doubly linked
list gives $O(1)$ ordered traversal, a natural place to attach
per-direction entry orientation, and a per-floor \code{size} for the
``scene $x$ of $y$'' indicator; branch and vertical connections that are
\emph{not} part of the walking sequence are held separately in
\code{crossReferences} (opposite corridor, floor up/down, lift array, room
entry, return-to-corridor). This separation is the core data-structure
decision.

\subsection{Engine Construction}
\code{buildPanoramaEngine} runs four passes. \emph{Pass~1} groups scenes by
floor, sorts each group by stored sequence index, and appends nodes to that
floor's list; forward entry orientation is the target's own calibrated
$(\mathrm{yaw},\mathrm{pitch})$ and backward entry orientation is the
source's yaw $+180^{\circ}$ --- both read live from calibration, never from
a stored (possibly stale) edge value. \emph{Pass~2} converts each scene's
non-chain hotspots into \code{crossReferences} pointers (an
\code{elevator} hotspot appends to the \code{lift} array; \code{upstairs}
sets \code{floorUp}; etc.). \emph{Pass~3} is cross-floor stitching
(Algorithm~\ref{alg:crossfloor}). \emph{Pass~4} attaches each
hand-authored cross-floor sightline hotspot to exactly the node it was
authored on, skipping rows whose source node cannot be resolved.

A read-time reclassification moves the ground floor's calibrated scene~$5$
--- a standalone courtyard view --- into its own single-scene floor bucket
``Central Quadrangle,'' closing the gap in the ground-floor sequence around
it. The underlying scene object is untouched; only its floor grouping
changes. This is why the tour's floor sequence has six entries over five
physical floor directories.

\subsection{Cross-Floor Stitching}
\begin{algorithm}[!t]
\caption{Cross-Floor Panorama Stitching (Pass 3)}
\label{alg:crossfloor}
\begin{algorithmic}[1]
\Require engine $\mathcal{E}$; ordered floor sequence
$\Phi = \langle\text{Entrance},\text{Ground},\text{First},\text{Second},\text{Third},\text{Quadrangle}\rangle$
\Ensure each floor's tail linked to the next floor's head
\For{$i \gets 1$ \textbf{to} $|\Phi|-1$}
  \State $L \gets \mathcal{E}.\text{getFloor}(\Phi_i)$;\ \ $U \gets \mathcal{E}.\text{getFloor}(\Phi_{i+1})$
  \If{$L{.}\text{tail} = \varnothing \ \lor\ U{.}\text{head} = \varnothing$} \textbf{continue} \EndIf
  \State $a \gets L{.}\text{tail}$;\ \ $b \gets U{.}\text{head}$
  \State $a{.}\text{next} \gets (\,b,\ \mathrm{entryYaw}{=}b{.}\text{yaw},\ \mathrm{entryPitch}{=}b{.}\text{pitch}\,)$
  \State $b{.}\text{previous} \gets (\,a,\ \mathrm{entryYaw}{=}a{.}\text{yaw}{+}180^{\circ},\ \mathrm{entryPitch}{=}a{.}\text{pitch}\,)$
\EndFor
\end{algorithmic}
\end{algorithm}

Algorithm~\ref{alg:crossfloor} runs \emph{after} all six per-floor lists
exist. It wires a lower floor's tail to an upper floor's head using exactly
the same link shape and the same entry-orientation rule as an ordinary
same-floor link. Because manual \code{goToOffset($\pm1$)} and guided
auto-advance both read only \code{node.next}/\code{node.previous}, a scene
at the end of its floor transparently continues into the next floor's first
scene, \emph{with no code path aware that a boundary was crossed}. Each
floor's own list (\code{head}/\code{tail}/\code{size}) is untouched, so
``scene $x$ of $y$'' still counts within a floor. By design there is no
wraparound from the last floor back to the first.

\subsection{Manual Traversal}
\begin{algorithm}[!t]
\caption{Manual Tour Hotspot Derivation and Step}
\label{alg:tourwalk}
\begin{algorithmic}[1]
\Require current node $v$; requested offset $\Delta\in\{-1,+1\}$ or hotspot $k$
\Ensure next scene id and entry orientation
\State $H \gets \emptyset$ \Comment{navigable hotspots, derived, never stored}
\If{$v{.}\text{next} \ne \varnothing$}\ $H \mathrel{+}= (\text{forward},\ v{.}\text{next})$ at yaw $v{.}\text{yaw}$ \EndIf
\If{$v{.}\text{previous} \ne \varnothing$}\ $H \mathrel{+}= (\text{back},\ v{.}\text{previous})$ at yaw $v{.}\text{yaw}{+}180^{\circ}$ \EndIf
\ForAll{$c \in \{$oppositeCorridor, floorUp, floorDown, lift$[\,]$, roomEntry, returnTo$\}$}
  \If{$v{.}\text{crossReferences}[c] \ne \varnothing$ \textbf{and} a raw edge angle exists}
     \State $H \mathrel{+}= (c,\ \text{target})$ at that edge's $(\mathrm{yaw},\mathrm{pitch})$
  \EndIf
\EndFor
\ForAll{$s \in v{.}\text{crossFloorHotspots}$}\ $H \mathrel{+}= (\text{crossFloor},\ s.\text{target})$ at $(s.\text{yaw},s.\text{pitch})$ \EndFor
\If{step is $\Delta$}
  \State $\ell \gets (\Delta = +1)\ ?\ v{.}\text{next} : v{.}\text{previous}$;\ \textbf{return} $(\ell{.}\text{node.id},\ \ell{.}\text{entryYaw},\ \ell{.}\text{entryPitch})$
\Else
  \State \textbf{return} target of hotspot $k$ with its authored entry orientation
\EndIf
\end{algorithmic}
\end{algorithm}
Every navigable hotspot is derived from the node's own links and
cross-references on each render (Algorithm~\ref{alg:tourwalk}); there is no
stored hotspot array and no manual enable/disable switch. Forward/back
hotspots are placed at the node's calibrated yaw and yaw~$+180^{\circ}$;
junction hotspots use the real edge angle because a corridor branch has no
fixed offset from ``forward.'' A lift with more than one reachable floor
renders a floor-select submenu.

\subsection{Guided Tour}
The guided tour is an automatic node-to-node walk driven entirely by
\code{node.next} object references (Fig.~\ref{fig:guided},
Algorithm~\ref{alg:guided}). At each node a fixed six-phase sequence runs:
\emph{pause} ($900$\,ms) $\to$ \emph{look-left} ($900$\,ms, $-40^{\circ}$)
$\to$ \emph{recenter} ($800$\,ms) $\to$ \emph{look-right} ($900$\,ms,
$+40^{\circ}$) $\to$ \emph{recenter} ($800$\,ms) $\to$ \emph{move}
($400$\,ms), with all durations scaled by a speed factor
($\times1.6$ slow, $\times1$ normal, $\times0.55$ fast). Correctness relies
on two mechanisms: a monotonically increasing \emph{cancellation token}
that is incremented on every \code{stop}/\code{restart}/unmount so that any
in-flight timer continuation observes a token mismatch and returns; and a
per-run \emph{visited set} that aborts the walk to an \emph{error} state if
\code{node.next} ever points to an already-visited scene (a corrupt-graph
guard). Pause freezes at a phase boundary after the in-flight camera
animation completes; resume continues from exactly that phase index;
restart rewinds to the node the tour was seeded on. Advancing uses the same
\code{goToScene} path as a manual hotspot click, so the fade transition is
identical between modes.

\begin{algorithm}[!t]
\caption{Guided Tour Execution at a Node}
\label{alg:guided}
\begin{algorithmic}[1]
\Require node $v$; start phase index $j$; run token $\tau$
\Ensure camera walk; advance or terminate
\For{$i \gets j$ \textbf{to} $5$}
  \If{$\tau \ne \tau_{\text{current}}$}\ \textbf{return} \Comment{stopped/restarted} \EndIf
  \If{status $=$ \emph{paused}}\ save $i$;\ \textbf{return} \EndIf
  \State render phase $\phi_i$;\ \textbf{await} camera animation for
         $\lceil \text{base}(\phi_i)\cdot \text{scale(speed)}\rceil$ ms
\EndFor
\State $w \gets v{.}\text{next.node}$
\If{$w = \varnothing$}\ status $\gets$ \emph{completed};\ \textbf{return} \EndIf
\If{$\neg\textsc{Usable}(w)$ \textbf{or} $w{.}\text{id} \in \text{visited}$}\ status $\gets$ \emph{error};\ \textbf{return} \EndIf
\State visited $\gets$ visited $\cup\ \{w{.}\text{id}\}$
\State \textsc{Advance}$(w{.}\text{id},\ v{.}\text{next.entryYaw},\ v{.}\text{next.entryPitch})$
\State \textbf{recurse} with $(w, 0, \tau)$
\end{algorithmic}
\end{algorithm}

\subsection{Minimap Projection}
The minimap (Algorithm~\ref{alg:minimap}) is inline SVG. For the current
floor it collects scenes whose backing node has real
$(\text{pos\_x},\text{pos\_y})$ coordinates; if at least two exist it
fits their bounding box into an on-card plot region with an
aspect-preserving scale and draws the real edges between them. If fewer
than two positioned scenes exist it falls back to an evenly spaced strip
inside the same region --- it \emph{never invents a coordinate}. A live
heading needle reads the panorama viewer's yaw once per animation frame and
is rotated by direct DOM mutation, bypassing React state to avoid 60\,Hz
re-renders. The outer ``entrance zone / building zone'' frame is explicitly
decorative and not survey-derived.

\begin{algorithm}[!t]
\caption{Minimap Coordinate Projection}
\label{alg:minimap}
\begin{algorithmic}[1]
\Require current floor scenes $F$; node positions $\text{pos}[\cdot]$; plot rect $Z=(x,y,w,h)$; pad $p$
\Ensure screen points for markers and edges
\State $S \gets \{(s,\text{pos}[s]) : s\in F,\ \text{pos}[s]\ \text{defined}\}$
\If{$|S| \ge 2$}
  \State $(x_{\min},x_{\max},y_{\min},y_{\max}) \gets$ bounding box of $S$
  \State $\sigma \gets \min\!\big(\tfrac{w-2p}{\max(x_{\max}-x_{\min},1)},\ \tfrac{h-2p}{\max(y_{\max}-y_{\min},1)}\big)$
  \State center the scaled box in $Z$; \textbf{for each} $(s,(a,b))$:
         $c_x \gets o_x + (a-x_{\min})\sigma$, $c_y \gets o_y + (b-y_{\min})\sigma$
  \State draw edge $(u,v)$ iff both endpoints are in $S$
\Else
  \State place $|F|$ markers evenly on the horizontal midline of $Z$ \Comment{no coordinate invented}
\EndIf
\end{algorithmic}
\end{algorithm}
